\section{Inline Code}\label{inline-code}

In C the \texttt{strstr} function defined with the prototype
{\ttfamily \PY{k+kt}{char}\PY{+w}{ }\PY{o}{*}\PY{n+nf}{strstr}\PY{p}{(}\PY{k}{const}\PY{+w}{ }\PY{k+kt}{char}\PY{+w}{ }\PY{o}{*}\PY{n}{haystack}\PY{p}{,}\PY{+w}{ }\PY{k}{const}\PY{+w}{ }\PY{k+kt}{char}\PY{+w}{ }\PY{o}{*}\PY{n}{needle}\PY{p}{)}\PY{p}{;}
} is
used to locate a substring within a string. It returns a pointer to the first occurrence of the substring
\texttt{needle} in the string \texttt{haystack}, or \texttt{NULL} if the substring is not found.

In Python, you can achieve similar functionality using the \texttt{find} method of strings for example: {\ttfamily \PY{n}{haystack}\PY{o}{.}\PY{n}{find}\PY{p}{(}\PY{n}{sub}\PY{p}{:}\PY{n+nb}{int}\PY{p}{)}\PY{o}{\PYZhy{}}\PY{o}{\PYZgt{}}\PY{n+nb}{int}
}. This method returns the lowest index of the substring if found in the string, otherwise it returns \texttt{-\allowbreak{}1}.
